\documentclass{beamer}

\def\Tiny{\fontsize{6pt}{6pt}\selectfont}
\def\supertiny{\fontsize{4pt}{4pt}\selectfont}

\mode<presentation>
{
  \usetheme{Warsaw}
  % \setbeamercovered{transparent}
  \usecolortheme{crane}
}

\usepackage{graphicx, ifthen, listings, fancyvrb}

\usepackage[czech]{babel}
% \usefonttheme{professionalfonts}
\usepackage{times}
\usepackage{amsmath}
\usepackage[utf8]{inputenc}
\usepackage{wrapfig}

\usepackage[T1]{fontenc}

\lstset{ basicstyle=\tiny, stringstyle=\ttfamily, showstringspaces=false }

\everymath{\displaystyle}

\setbeamerfont{frametitle}{size=\large}
\setbeamerfont{subsection in toc}{size=\scriptsize}

\makeatletter\newenvironment{blackbox}{%
   \begin{lrbox}{\@tempboxa}\begin{minipage}{0.95\columnwidth}}{\end{minipage}\end{lrbox}%
   \colorbox{black}{\usebox{\@tempboxa}}
}\makeatother

\title[IMF (11)]{Informatika pro moderní fyziky (8)\\ web scraping, procvičení ERb+LaTeX, tvorba vektorové grafiky}

\author[Franti\v{s}ek HAVL\r{U}J, ORF ÚJV Řež]{Franti\v{s}ek HAVL\r{U}J\\{\scriptsize \emph{e-mail: haf@ujv.cz}}}

\date{akademický rok 2015/2016\\23. listopadu 2016}

\institute[ORF ÚJV Řež]
{ÚJV Řež\\oddělení Reaktorové fyziky a podpory palivového cyklu}

\AtBeginSection[]
{
\begin{frame}<beamer>
\frametitle{Obsah}
\tableofcontents[currentsection,hideothersubsections]
\end{frame}
}

\begin{document}

\begin{frame}
  \titlepage
\end{frame}

\begin{frame}
  \tableofcontents
\end{frame}

\section{Komiks (= web scraping)}

\begin{frame}{Zpracování cizích zdrojů na webu -- web scraping}
  \begin{itemize}
    \item dosud jsme zpracovávali pouze lokální, hezky formátovaná data
    \item i leckteré externí služby poskytují pěkná API ve formátech JSON nebo XML
    \item ale leckdy taky ne a zajímavé informace
    \item protože jsme chytré horákyně, naučíme se, jak data automaticky získat
  \end{itemize}
\end{frame}

\begin{frame}{Zadání úkolu}
  \begin{itemize}
    \item protože po práci si chceme oddychnout a netrápit se přitom náročnou intelektuální činností, přečteme si rádi dobrý komiks
    \item a protože jsme přiměřeně cyničtí a také si chceme pocvičit angličtinu, přečteme si redmeat
    \item www.xkcd.com
    \item .. ale nechceme klikat a nechceme číst na internetu, takže si zhotovíme PDFko se všemi díly najednou
    \item zvládneme to snadno v LaTeXu, ale potřebujeme postahovat ty obrázky
  \end{itemize}
\end{frame}

\begin{frame}{HTML, }
  \begin{itemize}
    \item http://xkcd.com/1000
    \item zobrazíme zdrojový kód
    \item HTML: tagy, atributy, \texttt{class}, \texttt{id}
    \item naštěstí to první zajímavé umíme vykoukat bez znalostí... stačí regex
  \end{itemize}
\end{frame}

\begin{frame}[fragile]{Stahování webové stránky}
  Součást standardní knihovny -- \texttt{open-uri}
  \scriptsize
\begin{verbatim}
  require 'open-uri'
  f = open(remote_url)
  s = f.read
  m = s.match(...)
\end{verbatim}
\end{frame}

\begin{frame}{}
  \begin{itemize}
    \item najdu si url obrázku
    \item někam si ho uložím
    \item budu chytrý ohledně jména souboru
    \item dám bacha na příponu!
  \end{itemize}
\end{frame}

\begin{frame}[fragile]{Stahování dat}
  \scriptsize
\begin{verbatim}
  require 'open-uri'
  File.open(local_filename, 'wb') do |f2|
    open(remote_url, 'rb') do |f1|
      f2.write f1.read
    end
  end
\end{verbatim}
(kdo chce mít lepší život, tak si samozřejmě nadefinuje funkci!)
\end{frame}

\begin{frame}{Vygenerovat PDF}
  \begin{itemize}
    \item zase triviální, už to umíme, rychlá akce na deset minut!
    \item ERb šablona, použít \texttt{erb\_compiler} (mám z minula)
    \item seznam souborů vzít z \texttt{Dir["comics\_*"]}
    \item všechno známe z minula
  \end{itemize}
\end{frame}

\begin{frame}[fragile]{Další krok: popisky, nadpis a transkript}
  \begin{itemize}
    \item u každého komiksu je popisek -- atribut \texttt{title} u \texttt{img} tagu
    \item správný selektor je \texttt{\#comics img}
    \item nainstalujeme gem \texttt{nokogiri}
  \end{itemize}
  {\scriptsize
  \begin{verbatim}
f = open("http://xkcd.com/1000")
s = f.read
doc = Nokogiri::HTML(s)
img = doc.css("#comics img").first
puts img.attributes["title"]
  \end{verbatim}
  }
\end{frame}

\begin{frame}{}
  \begin{itemize}
    \item je tam navíc název a transkript
    \item je chytré to uložit chytře! takže ne 3 pole \texttt{titles}, \texttt{comments}, \texttt{transcripts} ale hezky jedno pole hashů
    \item zde přijdou ke slovu symboly: je to lepší klíč než řetězec
    \item \texttt{\{:title => ``...'', :comment => ``...''\}}
    \item  alternativní zápis
    \item \texttt{\{title: ``...'', comment: ``...''\}}
    \item a opět poskládám v PDF
  \end{itemize}
\end{frame}

\section{Tvorba obrázků}

\begin{frame}{Zadání dnešní úlohy}
  \begin{itemize}
    \item pro zadanou textovou mapu AZ VR1 potřebuju udělat hezký obrázek
    \item co druh, to barvička, rozumně zacházet s odstíny (palivo různě modré, R/B/E tyče různě červené, zelené, fialové)
  \end{itemize}
\end{frame}

\begin{frame}{Jak na obrázky}
  \begin{itemize}
    \item pěkný formát na tvorbu vektorových obrázku je SVG (Scalable Vector Graphics)
    \item je to dobrá věc především na internet -- všechny prohlížeče ho umí
    \item stejně jako HTML je postaven na XML
  \end{itemize}
\end{frame}

\begin{frame}[fragile]{Jednoduchý příklad}
  \tiny
  \begin{verbatim}
    <svg width="320" height="320" xmlns="http://www.w3.org/2000/svg" version="1.1">
      <rect x="0.0" y="0.0" width="40.0" height="40.0" fill="blue" />
      <rect x="40.0" y="0.0" width="40.0" height="40.0" fill="red" />
      <rect x="0.0" y="40.0" width="40.0" height="40.0" fill="green" />
      <rect x="40.0" y="40.0" width="40.0" height="40.0" fill="yellow" />
    </svg>
  \end{verbatim}
  \includegraphics[width=0.15\textwidth]{example}
\end{frame}

\begin{frame}{SVG -- co a jak}
  \begin{itemize}
    \item souřadný systém z levého horního rohu
    \item je potřeba udat celkovou šířku a výšku
    \item zatím nám stačí obdélník -- tag \texttt{rect}
    \item pozor, je to striktní XML, tedy je nutné \texttt{rect} tag uzavřít (!)
    \item vyzkoušejte -- nejdřív jen tak, potom vygenerovat 8x8 mapu (zatím klidně prázdnou)
  \end{itemize}
\end{frame}

\begin{frame}{Další SVG chytrosti}
  \begin{itemize}
    \item kromě \texttt{rect} se bude hodit také \texttt{text}
    \item jako text se zobrazí obsah příslušného elementu
    \item opět použiju atributy \texttt{x}, \texttt{y} (levý dolní roh) a můžu přihodit \texttt{text-anchor=``middle''}, aby to byl dolní prostředek
  \end{itemize}
\end{frame}

\begin{frame}{Postup}
  \begin{itemize}
    \item načtu ze souboru třeba do 2D pole
    \item budu mít hash s barvičkama
    \item vykreslím do SVG
  \end{itemize}
\end{frame}

\begin{frame}{A to je vše, přátelé!}
  \begin{center}
    \includegraphics[width=0.8\textwidth]{looney_tunes}
  \end{center}
\end{frame}

\end{document}
